
\documentclass[11pt,a4paper]{article}
\usepackage[T1]{fontenc}
\usepackage{lmodern}
\usepackage{amsmath,amssymb,amsfonts,mathtools}
\usepackage{graphicx,booktabs,longtable,array,tabularx,adjustbox}
\usepackage{geometry,microtype,fancyhdr,lastpage,etoolbox,enumitem}
\usepackage{newunicodechar,textcomp}
\usepackage{hyperref}
\hypersetup{hidelinks,pdfcreator={XeLaTeX / QM1 academic edition}}
\geometry{left=20mm,right=20mm,top=22mm,bottom=23mm}
\setcounter{secnumdepth}{2}
\setlength{\parindent}{0pt}
\setlength{\parskip}{4pt plus 1pt}
\setlength{\emergencystretch}{2em}
\setlength{\tabcolsep}{3pt}
\renewcommand{\arraystretch}{1.16}
\setlist{nosep,topsep=3pt,itemsep=2pt,leftmargin=7mm}
\AtBeginEnvironment{longtable}{\footnotesize\setlength{\tabcolsep}{2.2pt}}
\newcommand{\tightlist}{\setlength{\itemsep}{0pt}\setlength{\parskip}{0pt}}
\providecommand{\Cov}{\operatorname{Cov}}
\providecommand{\Var}{\operatorname{Var}}
\providecommand{\Corr}{\operatorname{Corr}}
\providecommand{\MAD}{\operatorname{MAD}}
\providecommand{\E}{\mathbb{E}}
\providecommand{\R}{\mathbb{R}}
\newunicodechar{—}{\textemdash}
\newunicodechar{–}{\textendash}
\newunicodechar{€}{\texteuro}
\newunicodechar{≤}{\ensuremath{\leq}}
\newunicodechar{≥}{\ensuremath{\geq}}
\newunicodechar{×}{\ensuremath{\times}}
\newunicodechar{÷}{\ensuremath{\div}}
\newunicodechar{Σ}{\ensuremath{\Sigma}}
\newunicodechar{σ}{\ensuremath{\sigma}}
\newunicodechar{μ}{\ensuremath{\mu}}
\newunicodechar{Δ}{\ensuremath{\Delta}}
\newunicodechar{∑}{\ensuremath{\sum}}
\newunicodechar{≈}{\ensuremath{\approx}}
\newunicodechar{√}{\ensuremath{\surd}}
\newunicodechar{≠}{\ensuremath{\ne}}
\newunicodechar{→}{\ensuremath{\to}}
\newunicodechar{∞}{\ensuremath{\infty}}
\newunicodechar{∈}{\ensuremath{\in}}
\newunicodechar{−}{\ensuremath{-}}
\newunicodechar{²}{\ensuremath{^{2}}}
\newunicodechar{³}{\ensuremath{^{3}}}
\newunicodechar{₀}{\ensuremath{_{0}}}
\newunicodechar{₁}{\ensuremath{_{1}}}
\newunicodechar{₂}{\ensuremath{_{2}}}
\newunicodechar{₃}{\ensuremath{_{3}}}
\newunicodechar{₄}{\ensuremath{_{4}}}
\newunicodechar{ₜ}{\ensuremath{_{t}}}
\newunicodechar{ₚ}{\ensuremath{_{p}}}
\newunicodechar{ₙ}{\ensuremath{_{n}}}
\newunicodechar{ₓ}{\ensuremath{_{x}}}
\newunicodechar{ᵢ}{\ensuremath{_{i}}}
\newunicodechar{ŷ}{\ensuremath{\hat y}}
\newunicodechar{Ŷ}{\ensuremath{\widehat Y}}
\newunicodechar{ₘ}{\ensuremath{_{m}}}
\newunicodechar{ₑ}{\ensuremath{_{e}}}
\newunicodechar{ₐ}{\ensuremath{_{a}}}
\newunicodechar{ₖ}{\ensuremath{_{k}}}
\newunicodechar{ₛ}{\ensuremath{_{s}}}
\newunicodechar{ₗ}{\ensuremath{_{l}}}
\newunicodechar{ₕ}{\ensuremath{_{h}}}
\newunicodechar{ₒ}{\ensuremath{_{o}}}
\newunicodechar{₊}{\ensuremath{_{+}}}
\newunicodechar{₋}{\ensuremath{_{-}}}
\newunicodechar{₌}{\ensuremath{_{=}}}
\newunicodechar{₍}{\ensuremath{_{(}}}
\newunicodechar{₎}{\ensuremath{_{)}}}
\newunicodechar{⁻}{\ensuremath{^{-}}}
\newunicodechar{⁺}{\ensuremath{^{+}}}
\newunicodechar{⁰}{\ensuremath{^{0}}}
\newunicodechar{¹}{\ensuremath{^{1}}}
\newunicodechar{⁴}{\ensuremath{^{4}}}
\newunicodechar{⁵}{\ensuremath{^{5}}}
\newunicodechar{⁶}{\ensuremath{^{6}}}
\newunicodechar{⁷}{\ensuremath{^{7}}}
\newunicodechar{⁸}{\ensuremath{^{8}}}
\newunicodechar{⁹}{\ensuremath{^{9}}}
\makeatletter
\renewcommand{\section}{\@startsection{section}{1}{0pt}{2.0ex plus .7ex minus .3ex}{1.0ex}{\large\bfseries}}
\renewcommand{\subsection}{\@startsection{subsection}{2}{0pt}{1.8ex plus .5ex minus .2ex}{.7ex}{\normalsize\bfseries}}
\renewcommand{\subsubsection}{\@startsection{subsubsection}{3}{0pt}{1.3ex plus .4ex minus .2ex}{.5ex}{\normalsize\bfseries}}
\makeatother
\pagestyle{fancy}
\fancyhf{}
\renewcommand{\headrulewidth}{0.35pt}
\renewcommand{\footrulewidth}{0pt}
\fancyhead[L]{\small\textsc{KEDGE Business School --- QM1}}
\fancyhead[R]{\small\textsc{2026--2027}}
\fancyfoot[L]{\small Badr ABOUFARASSE, PhD}
\fancyfoot[R]{\small\thepage\ /\ \pageref{LastPage}}
\setlength{\headheight}{14pt}
\setkeys{Gin}{keepaspectratio}
\begin{document}

\begin{center}
{\small\textsc{KEDGE Business School \quad EBP1 \quad Quantitative Methods 1}}\\[6pt]
{\Large\bfseries Fiche 5 — Séries temporelles}\\[6pt]
{\normalsize\bfseries Corrigés intégralement détaillés}\\[6pt]
{\small Badr ABOUFARASSE, PhD \quad --- \quad 2026--2027}
\end{center}
\vspace{2mm}\hrule\vspace{6mm}

\hypertarget{applications-fondamentales}{%
\section{Applications fondamentales}\label{applications-fondamentales}}

Maîtriser les définitions et les calculs essentiels du cours.

\emph{Application du cours --- 10 min}

\hypertarget{exercice-1-composantes-et-moduxe8les}{%
\subsection{Exercice 1 -- Composantes et modèles}\label{exercice-1-composantes-et-moduxe8les}}

\begin{enumerate}
\item
  Associer à {\(T_t\)}, {\(C_t\)}, {\(S_t\)} ou {\(R_t\)} : tendance de long terme, récession, pic récurrent du quatrième trimestre, choc exceptionnel.
\item
  Écrire le modèle additif complet et le modèle multiplicatif complet.
\item
  Expliquer dans quel cas une saisonnalité multiplicative est plus plausible qu'une saisonnalité additive.
\end{enumerate}

\hypertarget{corriguxe9-duxe9tailluxe9}{%
\subsubsection{Corrigé détaillé}\label{corriguxe9-duxe9tailluxe9}}

{\[\begin{gathered}
T_t:\text{ tendance de long terme},\\C_t:\text{ cycle},\\S_t:\text{ saisonnalité},\\R_t:\text{ irrégulier/résiduel}.
\end{gathered}\]} Modèle additif : {\[Y_t=T_t+C_t+S_t+R_t.\]} Modèle multiplicatif : {\[Y_t=T_t\times C_t\times
S_t\times R_t.\]} Le modèle multiplicatif est particulièrement plausible lorsque l'amplitude saisonnière augmente avec le niveau de la série.

\hypertarget{distinguer-les-unituxe9s-et-les-facteurs}{%
\subsubsection{Distinguer les unités et les facteurs}\label{distinguer-les-unituxe9s-et-les-facteurs}}

Une composante saisonnière additive s'exprime en unités de Y et conserve la même amplitude absolue. Une composante saisonnière multiplicative s'exprime comme un facteur (par exemple 1,20 pour +20 \%), dont l'effet absolu varie avec le niveau de tendance. Le cycle correspond à des fluctuations de plus longue durée, tandis qu'un choc irrégulier n'est pas périodique.

\emph{Application du cours --- 20 min}

\hypertarget{exercice-2-moyennes-mobiles-et-centrage}{%
\subsection{Exercice 2 -- Moyennes mobiles et centrage}\label{exercice-2-moyennes-mobiles-et-centrage}}

On observe deux années de demande trimestrielle : {\[80,\ 100,\ 120,\ 140,\ 88,\ 110,\ 132,\
154.\]}

\begin{enumerate}
\item
  Décrire le motif principal de la série.
\item
  Calculer les cinq moyennes mobiles d'ordre 4.
\item
  Les centrer afin d'obtenir les moyennes mobiles centrées aux périodes 3 à 6.
\item
  Expliquer pourquoi une moyenne mobile d'ordre pair doit être recentrée.
\end{enumerate}

\hypertarget{corriguxe9-duxe9tailluxe9-1}{%
\subsubsection{Corrigé détaillé}\label{corriguxe9-duxe9tailluxe9-1}}

{\[MA_{2.5}=\frac{80+100+120+140}{4}=110.\]} {\[MA_{3.5}=\frac{100+120+140+88}{4}=112.\]} {\[MA_{4.5}=\frac{120+140+88+110}{4}=114.5.\]} {\[MA_{5.5}=\frac{140+88+110+132}{4}=117.5.\]} {\[MA_{6.5}=\frac{88+110+132+154}{4}=121.\]} Les moyennes mobiles centrées sont : {\[CMA_3=\frac{110+112}{2}=111,\]} {\[CMA_4=\frac{112+114.5}{2}=113.25,\]} {\[CMA_5=\frac{114.5+117.5}{2}=116,\]} {\[CMA_6=\frac{117.5+121}{2}=119.25.\]} Une moyenne d'ordre 4 se situe entre deux périodes. Le centrage replace l'estimation sur une période observée.

\hypertarget{tableau-dalignement-des-moyennes-mobiles}{%
\subsubsection{Tableau d'alignement des moyennes mobiles}\label{tableau-dalignement-des-moyennes-mobiles}}

\begin{longtable}[]{@{}llll@{}}
\toprule\noalign{}
Fenêtre & Somme & MA4 & Centre \\
\midrule\noalign{}
\endhead
\bottomrule\noalign{}
\endlastfoot
1--4 & 440 & 110 & 2.5 \\
2--5 & 448 & 112 & 3.5 \\
3--6 & 458 & 114.5 & 4.5 \\
4--7 & 470 & 117.5 & 5.5 \\
5--8 & 484 & 121 & 6.5 \\
\end{longtable}

\begin{longtable}[]{@{}lllll@{}}
\toprule\noalign{}
t & 3 & 4 & 5 & 6 \\
\midrule\noalign{}
\endhead
\bottomrule\noalign{}
\endlastfoot
CMA4 & (110+112)/2=111 & (112+114.5)/2=113.25 & (114.5+117.5)/2=116 & (117.5+121)/2=119.25 \\
\end{longtable}

L'ordre 4 étant pair, la première moyenne mobile est centrée entre t=2 et t=3. Le centrage fournit des valeurs aux dates entières t=3 à 6.

\hypertarget{pruxe9paration-aux-uxe9valuations}{%
\section{Préparation aux évaluations}\label{pruxe9paration-aux-uxe9valuations}}

Résoudre des exercices de contrôle continu et de partiel en autonomie.

\emph{Préparation CC / partiel --- 25 min}

\hypertarget{exercice-3-duxe9composition-multiplicative-et-pruxe9vision}{%
\subsection{Exercice 3 -- Décomposition multiplicative et prévision}\label{exercice-3-duxe9composition-multiplicative-et-pruxe9vision}}

\begin{longtable}[]{@{}lllll@{}}
\toprule\noalign{}
Année & Q1 & Q2 & Q3 & Q4 \\
\midrule\noalign{}
\endhead
\bottomrule\noalign{}
\endlastfoot
2005 & 82 & 98 & 120 & 150 \\
2006 & 91 & 111 & 137 & 171 \\
2007 & 104 & 126 & 154 & 194 \\
\end{longtable}

On donne la tendance estimée : {\[\widehat
T_t=6.343t+86.939.\]} Les coefficients saisonniers normalisés sont : {\[S_{Q1}=0.789,\quad S_{Q2}=0.903,\quad
S_{Q3}=1.052,\quad S_{Q4}=1.255.\]}

\begin{enumerate}
\item
  Interpréter la pente de la tendance.
\item
  Vérifier approximativement la normalisation des coefficients saisonniers.
\item
  Calculer les quatre prévisions trimestrielles pour 2008.
\item
  Expliquer comment désaisonnaliser une observation {\(Y_t\)} dans le modèle multiplicatif.
\end{enumerate}

\hypertarget{corriguxe9-duxe9tailluxe9-2}{%
\subsubsection{Corrigé détaillé}\label{corriguxe9-duxe9tailluxe9-2}}

La pente 6,343 signifie une hausse tendancielle de 6,343 unités par trimestre. Les coefficients saisonniers fournis totalisent 3,999, soit environ 4 (écart d\textquotesingle arrondi).

\begin{longtable}[]{@{}lllll@{}}
\toprule\noalign{}
t & Q & Tₜ=6.343t+86.939 & S\textquotesingle{} & Tₜ×S\textquotesingle{} \\
\midrule\noalign{}
\endhead
\bottomrule\noalign{}
\endlastfoot
13 & Q1 & 169.398 & 0.789 & 133.655022 \\
14 & Q2 & 175.741 & 0.903 & 158.694123 \\
15 & Q3 & 182.084 & 1.052 & 191.552368 \\
16 & Q4 & 188.427 & 1.255 & 236.475885 \\
\end{longtable}

\[\widehat Y_{13}=169.398(0.789)=133.655022,\quad\widehat Y_{14}=175.741(0.903)=158.694123.\]

\[\widehat Y_{15}=182.084(1.052)=191.552368,\quad\widehat Y_{16}=188.427(1.255)=236.475885.\]

Pour désaisonnaliser une observation, diviser Yₜ par le coefficient S\textquotesingle{} du trimestre correspondant. Garder les valeurs non arrondies jusqu\textquotesingle au résultat final.

\emph{Préparation CC / partiel --- 65 min}

\hypertarget{exercice-5-chauxeene-compluxe8te-de-duxe9composition-saisonniuxe8re-type-partiel}{%
\subsection{Exercice 5 --- Chaîne complète de décomposition saisonnière --- type partiel}\label{exercice-5-chauxeene-compluxe8te-de-duxe9composition-saisonniuxe8re-type-partiel}}

Une entreprise suit ses ventes trimestrielles pendant trois ans. Utiliser le modèle multiplicatif Yₜ=Tₜ×Sₜ, sans cycle ni composante irrégulière explicites. Ne pas arrondir les résultats intermédiaires.

\begin{longtable}[]{@{}lllll@{}}
\toprule\noalign{}
Année / Year & Q1 & Q2 & Q3 & Q4 \\
\midrule\noalign{}
\endhead
\bottomrule\noalign{}
\endlastfoot
2019 & 40 & 52 & 48 & 68 \\
2020 & 44 & 58 & 53 & 76 \\
2021 & 50 & 65 & 60 & 84 \\
\end{longtable}

\begin{enumerate}
\tightlist
\item
  Tracer la série et calculer les neuf moyennes mobiles d\textquotesingle ordre 4.
\item
  Centrer les moyennes mobiles (t=3 à 10). Construire le tableau t, CMA4, t², t×CMA4.
\item
  Estimer la tendance Tₜ=at+b par moindres carrés sur les huit CMA4.
\item
  Calculer Yₜ/Tₜ pour les 12 trimestres, les quatre moyennes saisonnières et leurs coefficients rectifiés de somme 4.
\item
  Prévoir les quatre trimestres de 2022 (t=13 à 16), en valeurs désaisonnalisées puis saisonnalisées.
\end{enumerate}

\hypertarget{corriguxe9-duxe9tailluxe9-3}{%
\subsubsection{Corrigé détaillé}\label{corriguxe9-duxe9tailluxe9-3}}

1. Moyennes mobiles d\textquotesingle ordre 4 :

\[MA_{2.5}=(40+52+48+68)/4=52,\quad MA_{3.5}=(52+48+68+44)/4=53.\]

\begin{longtable}[]{@{}lllll@{}}
\toprule\noalign{}
Fenêtre / Window & Périodes / Periods & Somme / Sum & MA4 & Centre / Midpoint \\
\midrule\noalign{}
\endhead
\bottomrule\noalign{}
\endlastfoot
1 & 1--4 & 208 & 52 & 2.5 \\
2 & 2--5 & 212 & 53 & 3.5 \\
3 & 3--6 & 218 & 54.5 & 4.5 \\
4 & 4--7 & 223 & 55.75 & 5.5 \\
5 & 5--8 & 231 & 57.75 & 6.5 \\
6 & 6--9 & 237 & 59.25 & 7.5 \\
7 & 7--10 & 244 & 61 & 8.5 \\
8 & 8--11 & 251 & 62.75 & 9.5 \\
9 & 9--12 & 259 & 64.75 & 10.5 \\
\end{longtable}

2. Centrage : les neuf MA4, situées à des demi-périodes, donnent huit CMA4 aux périodes entières 3 à 10.

\[CMA_3=(52+53)/2=52.5,\quad CMA_4=(53+54.5)/2=53.75.\]

\begin{longtable}[]{@{}llll@{}}
\toprule\noalign{}
t & CMA4 & t² & t×CMA4 \\
\midrule\noalign{}
\endhead
\bottomrule\noalign{}
\endlastfoot
3 & 52.5 & 9 & 157.5 \\
4 & 53.75 & 16 & 215 \\
5 & 55.125 & 25 & 275.625 \\
6 & 56.75 & 36 & 340.5 \\
7 & 58.5 & 49 & 409.5 \\
8 & 60.125 & 64 & 481 \\
9 & 61.875 & 81 & 556.875 \\
10 & 63.75 & 100 & 637.5 \\
Σt=52 & 462.375 & 380 & 3073.5 \\
\end{longtable}

3. Estimation de la tendance à partir des CMA4 :

\[\bar t=52/8=6.5,\quad\overline{CMA}=462.375/8=57.796875.\]

\[V(t)=380/8-6.5^2=5.25,\quad\operatorname{Cov}(t,CMA)=3073.5/8-6.5(57.796875)=8.5078125.\]

\[a=8.5078125/5.25=1.620535714,\quad b=57.796875-a(6.5)=47.263392857.\]

\[\widehat T_t=1.620535714t+47.263392857.\]

4. Ratios saisonniers de chaque observation et normalisation :

\begin{longtable}[]{@{}lllll@{}}
\toprule\noalign{}
t & Trimestre / Quarter & Y & Tₜ & Sₜ=Y/Tₜ \\
\midrule\noalign{}
\endhead
\bottomrule\noalign{}
\endlastfoot
1 & Q1 & 40 & 48.8839 & 0.81826 \\
2 & Q2 & 52 & 50.5045 & 1.02961 \\
3 & Q3 & 48 & 52.125 & 0.92086 \\
4 & Q4 & 68 & 53.7455 & 1.26522 \\
5 & Q1 & 44 & 55.3661 & 0.79471 \\
6 & Q2 & 58 & 56.9866 & 1.01778 \\
7 & Q3 & 53 & 58.6071 & 0.90433 \\
8 & Q4 & 76 & 60.2277 & 1.26188 \\
9 & Q1 & 50 & 61.8482 & 0.80843 \\
10 & Q2 & 65 & 63.4688 & 1.02413 \\
11 & Q3 & 60 & 65.0893 & 0.92181 \\
12 & Q4 & 84 & 66.7098 & 1.25918 \\
\end{longtable}

\[\bar S_q=\frac{1}{3}\sum_{j=0}^2(Y_{4j+q}/\widehat T_{4j+q}),\qquad S'_q=\frac{4\bar S_q}{\sum_{k=1}^4\bar S_k}.\]

\begin{longtable}[]{@{}llll@{}}
\toprule\noalign{}
Trimestre / Quarter & Ratios / Ratios & Moyenne / Mean & S\textquotesingle{} normalisé / Normalised \\
\midrule\noalign{}
\endhead
\bottomrule\noalign{}
\endlastfoot
Q1 & 0.81826 ; 0.79471 ; 0.80843 & 0.807135 & 0.805376 \\
Q2 & 1.02961 ; 1.01778 ; 1.02413 & 1.02384 & 1.021609 \\
Q3 & 0.92086 ; 0.90433 ; 0.92181 & 0.915667 & 0.913671 \\
Q4 & 1.26522 ; 1.26188 ; 1.25918 & 1.262095 & 1.259344 \\
TOTAL & --- & 4.008737 & 4 \\
\end{longtable}

5. Prévision : Tₜ représente la prévision désaisonnalisée ; la prévision finale est Tₜ×S\textquotesingle{} du trimestre.

\begin{longtable}[]{@{}lllll@{}}
\toprule\noalign{}
t & Trimestre / Quarter & Tₜ & S\textquotesingle{} & Prévision / Forecast \\
\midrule\noalign{}
\endhead
\bottomrule\noalign{}
\endlastfoot
13 & Q1 & 68.3304 & 0.805376 & 55.0316 \\
14 & Q2 & 69.9509 & 1.021609 & 71.4624 \\
15 & Q3 & 71.5714 & 0.913671 & 65.3927 \\
16 & Q4 & 73.192 & 1.259344 & 92.1739 \\
\end{longtable}

\includegraphics[width=0.85\textwidth,height=\textheight]{assets/chapter-5-solutions-fr-chart-01.pdf}

\emph{Ventes observées et prévisions 2022 / Observed and 2022 forecasts}

Contrôle : la somme des quatre coefficients normalisés est égale à 4. Les prévisions finales 2022 sont environ 55,03 ; 71,46 ; 65,39 ; 92,17.

\hypertarget{approfondissement-facultatif}{%
\section{Approfondissement facultatif}\label{approfondissement-facultatif}}

Aller plus loin dans l'interprétation et les applications professionnelles.

\emph{Approfondissement --- 15 min}

\hypertarget{exercice-4-uxe9valuation-des-pruxe9visions}{%
\subsection{Exercice 4 -- Évaluation des prévisions}\label{exercice-4-uxe9valuation-des-pruxe9visions}}

Pour cet exercice, on suppose que les valeurs effectivement observées en 2008 ont été : {\[136,\ 162,\ 188,\
245.\]} À partir des prévisions de l'exercice 3 :

\begin{enumerate}
\item
  calculer les erreurs {\(e_t=Y_t-\hat Y_t\)} ;
\item
  calculer le MAE ;
\item
  calculer le RMSE ;
\item
  expliquer pourquoi une période de validation distincte de la période d'estimation est utile ;
\item
  rappeler pourquoi une prévision naïve constitue un benchmark utile.
\end{enumerate}

\hypertarget{corriguxe9-duxe9tailluxe9-4}{%
\subsubsection{Corrigé détaillé}\label{corriguxe9-duxe9tailluxe9-4}}

Les prévisions sont celles effectivement obtenues à l\textquotesingle exercice 3, sans substitution de chiffres arrondis issus d\textquotesingle un autre support.

\begin{longtable}[]{@{}llllll@{}}
\toprule\noalign{}
Q & Y & Prévision / Forecast & e=Y−Ŷ & \textbar e\textbar{} & e² \\
\midrule\noalign{}
\endhead
\bottomrule\noalign{}
\endlastfoot
1 & 136 & 133.655022 & 2.344978 & 2.344978 & 5.49892 \\
2 & 162 & 158.694123 & 3.305877 & 3.305877 & 10.92882 \\
3 & 188 & 191.552368 & -3.552368 & 3.552368 & 12.61932 \\
4 & 245 & 236.475885 & 8.524115 & 8.524115 & 72.66054 \\
\end{longtable}

\[MAE=(2.344978+3.305877+3.552368+8.524115)/4=4.431835.\]

\[RMSE=\sqrt{(2.344978^2+3.305877^2+(-3.552368)^2+8.524115^2)/4}=5.042509.\]

\[ME=(2.344978+3.305877-3.552368+8.524115)/4=2.655651.\]

Le modèle sous-prévoit en moyenne. La validation doit être hors échantillon ; une prévision naïve est un benchmark utile.

\emph{Approfondissement --- 20 min}

\hypertarget{compluxe9ment-1-validation-des-pruxe9visions-et-rupture-structurelle}{%
\subsection{Complément 1 --- Validation des prévisions et rupture structurelle}\label{compluxe9ment-1-validation-des-pruxe9visions-et-rupture-structurelle}}

La demande mensuelle observée vaut 100, 104, 108, 150, 154 et 158. Un modèle ajusté sur les trois premiers mois prévoit 100, 104, 108, 112, 116 et 120.

\begin{enumerate}
\tightlist
\item
  Calculez les erreurs des mois 4 à 6.
\item
  Calculez la MAE et l'erreur moyenne sur la validation.
\item
  Diagnostiquez le modèle.
\item
  Proposez une réponse responsable.
\end{enumerate}

\textbf{Indication méthodologique :}

Utilisez erreur = observé − prévu. Des erreurs toujours positives signalent une sous-prévision systématique.

\hypertarget{corriguxe9-duxe9tailluxe9-5}{%
\subsubsection{Corrigé détaillé}\label{corriguxe9-duxe9tailluxe9-5}}

\textbf{Erreurs.} Mois 4 : 150−112 = 38 ; mois 5 : 154−116 = 38 ; mois 6 : 158−120 = 38.

\textbf{Indicateurs.} MAE = \textbf{38} et erreur moyenne = \textbf{+38} : la sous-prévision est systématique.

\textbf{Diagnostic.} Le niveau saute entre les mois 3 et 4, puis la pente reste proche de quatre unités par mois. Il s'agit vraisemblablement d'une rupture de niveau ; le modèle pré-rupture est obsolète.

\textbf{Réponse.} Identifier l'événement économique, réestimer le niveau sur les données post-rupture, comparer à une prévision naïve, conserver une période de validation et communiquer des intervalles de prévision.

\end{document}
